\chapter{Custom SI Units \& Qualifiers Reference}

This appendix documents the additional SI units and unit qualifiers
defined by OmniLaTeX in \texttt{omnilatex-math.sty}.  These extend the
standard \texttt{siunitx} v3 unit set for domain-specific use in
engineering and the physical sciences.

\section{Custom SI Units}

All custom units are declared with
\texttt{\textbackslash DeclareSIUnit} and are available immediately
after module load.  They integrate with the standard
\texttt{\textbackslash unit}, \texttt{\textbackslash qty}, and
\texttt{\textbackslash qtyrange} commands.

\begin{table}[htbp]
  \centering
  \caption{Custom SI units defined by OmniLaTeX.}
  \label{tab:ref-custom-units}
  \begin{tabular}{@{} l l l p{5cm} @{}}
    \toprule
    \textbf{Command} & \textbf{Symbol} & \textbf{Base} &
    \textbf{Description} \\
    \midrule
    \texttt{\textbackslash volpercent}
      & vol\,\% & percent
      & Volume percent; declared as text to preserve hyphen \\[4pt]
    \texttt{\textbackslash watthour}
      & Wh & J/s $\times$ h
      & Watt-hour energy unit \\[4pt]
    \texttt{\textbackslash annum}
      & a & year
      & Annum (Julian year, $365.25$ days) \\[4pt]
    \texttt{\textbackslash atmosphere}
      & atm & pressure
      & Standard atmosphere ($101\,325$\,Pa) \\[4pt]
    \texttt{\textbackslash partspermillion}
      & ppm & ratio
      & Parts per million ($10^{-6}$) \\[4pt]
    \texttt{\textbackslash bar}
      & bar & pressure
      & Bar ($10^{5}$\,Pa) \\[4pt]
    \texttt{\textbackslash relhumidity}
      & \%rh / \%\,r.F. & ratio
      & Relative humidity; language-dependent label \\
    \bottomrule
  \end{tabular}
\end{table}

The \texttt{\textbackslash relhumidity} unit adapts to the document
language.  In German documents it renders as \texttt{\%\,r.F.}\
(\emph{relative Feuchte}); in English documents as \texttt{\%\,RH}.

\section{Custom SI Qualifiers}

Unit qualifiers appear as subscripts and disambiguate quantities that
share the same SI unit.  OmniLaTeX declares three qualifiers via
\texttt{\textbackslash DeclareSIQualifier}, sourcing their display text
from the glossary system for consistency.

\begin{table}[htbp]
  \centering
  \caption{Custom SI qualifiers defined by OmniLaTeX.}
  \label{tab:ref-si-qualifiers}
  \begin{tabular}{@{} l l p{6cm} @{}}
    \toprule
    \textbf{Command} & \textbf{Qualifier} & \textbf{Description} \\
    \midrule
    \texttt{\textbackslash dryair}
      & dry air
      & Dry air component in moist-air calculations;
        reads from glossary entries \texttt{sub.dry} and
        \texttt{sub.air} \\[4pt]
    \texttt{\textbackslash water}
      & water
      & Water component; reads from glossary entry
        \texttt{sub.water} \\[4pt]
    \texttt{\textbackslash thermal}
      & thermal
      & Thermal property qualifier; reads from glossary entry
        \texttt{sub.thermal} \\
    \bottomrule
  \end{tabular}
\end{table}

Usage follows the standard \texttt{siunitx} qualifier syntax:

\begin{verbatim}
\unit{\kilogram\per\cubic\metre\dryair}
% → kg/m³_dry air

\qty{1.2}{\kilo\watt\per\metre\thermal}
% → 1.2 kW/m_thermal
\end{verbatim}

The Euro sign is also available as a unit via the \texttt{eurosym}
package---use \texttt{\textbackslash euro} directly inside
\texttt{\textbackslash unit\{\}} or \texttt{\textbackslash qty\{\}\{}}.
